\documentclass[a4paper,10pt]{report}\input{../0.Préambule/Préambule_cours_prof}\begin{document}

\pagestyle{empty}
\newpage
\section*{Devoir surveillé 4 : Factorisation et équations}
\addcontentsline{toc}{section}{Devoir surveillé 4 : Factorisation et équations}

\noindent \textit{Il sera tenu compte dans la notation de la propreté ainsi que de la justification apportée à chacune des réponses.}

\noindent \textit{Le barème est donné à titre indicatif. Il pourra être modifié ultérieurement.}

\noindent \textit{L'usage de la calculatrice est \textbf{autorisé}.} \hfill \textit{Durée : 55 minutes}

\medskip
\paragraph{Exercice 1}:\quad Question de cours \hfill \textbf{2 points}

\begin{enumerate}
\item Cite la formule qui permet de factoriser à l'aide d'une \og identité remarquable \fg{}.
\item Cite la propriété qui permet de résoudre une équation a produit nul.
\end{enumerate}

\medskip
\paragraph{Exercice 2}: \quad d'après Brevet 2021 \hfill  \textbf{6 points}

\medskip
Voici quatre affirmations. Pour chacune d'entre elles, dire si elle est vraie ou fausse. On rappelle que la réponse doit être justifiée.

\medskip

\textbf{Affirmation 1 :} $\dfrac{1}{3}$ est solution de l'équation $4x+\dfrac{3}{4}= 1$\medskip
		
\textbf{Affirmation 2 :} pour tout nombre $x$, \quad $(x + 9)^2  = x^2+81$.\medskip

\textbf{Affirmation 3 :} L'équation $x^2 = 8$ a une unique solution.	\medskip

\textbf{Affirmation 4 :} pour tout nombre $x$, \quad $(2x + 1)^2 - 4 = (2x + 3)(2x - 1)$.\medskip


\medskip
\paragraph{Exercice 3} : \quad d'après Brevet 2021\hfill  \textbf{6 points}

\medskip
On considère le programme de calcul : \:\begin{tabular}{|l|}\hline
$\bullet~~$Choisir un nombre.\\
$\bullet~~$Prendre le carré de ce nombre.\\
$\bullet~~$Ajouter le triple du nombre de départ.\\
$\bullet~~$Ajouter 2.\\ \hline
\end{tabular}

\medskip

\begin{enumerate}
\item Montrer que si on choisit 1 comme nombre de départ, le programme donne $6$ comme résultat.
\item Quel résultat obtient-on si on choisit $-5$ comme nombre de départ ?
\item On appelle $x$ le nombre de départ, exprimer le résultat du programme en fonction de $x$.
\item En développant $(x + 2)(x + 1)$, montrer que ce résultat peut aussi s'écrire sous la forme $(x + 2)(x + 1)$ pour toutes les valeurs de $x$.
\item  La feuille du tableur suivante regroupe des résultats du programme de calcul précédent.

\medskip
\begin{tabularx}{\linewidth}{|c|l|*{9}{>{\centering \arraybackslash}X|}}\hline
	&A 			&B		&C		&D		&E		&F		&G		&H		&I		&J\\ \hline
1	&$x$		& $-4$	&$-3$	&$-2$	&$-1$	&0		&1		&2		&3		&4\\ \hline
2	&$(x+2)(x+1)$&6		&2 		&0 		&0 		&2 	&6 &12 &20 &30\\ \hline
\end{tabularx}
\medskip

	\begin{enumerate}
		\item Quelle formule a été écrite dans la cellule $B2$ avant de l'étendre jusqu'à la cellule $J2$ ?
		\item Trouver les valeurs de $x$ pour lesquelles le programme donne $0$ comme résultat.
 	\end{enumerate}
\end{enumerate}	

\medskip
\paragraph{Exercice 4}\hfill  \textbf{3 points}

\medskip
\noindent Factoriser les expressions suivantes :
\begin{enumerate}
\item $A = 4x^2 + 12x + 16$
\item $N = (x - 4) ^2 - 36$
\item $E =(15x+7)(3-x)+(15x+7)(12x+5)$
\end{enumerate} 

\medskip
\paragraph{Exercice 5}\hfill  \textbf{3 points}

\medskip
\noindent Résoudre les équations suivantes :
\begin{enumerate}
\item $10x+1=0$
\item $2x-9=8x+3$
\item $(-15x+3)(3x+9)=0$
\end{enumerate} 

\end{document}